% !TeX program = lualatex
\documentclass{article}
\usepackage{kotex}
\setmainhangulfont{Noto Serif KR}[
  Renderer=HarfBuzz,
  BoldFont={Noto Serif KR},
  BoldFeatures={Renderer=HarfBuzz,FakeBold=2},
  ItalicFont={Noto Serif KR},
  ItalicFeatures={Renderer=HarfBuzz,FakeSlant=0.2},
  BoldItalicFont={Noto Serif KR},
  BoldItalicFeatures={Renderer=HarfBuzz,FakeBold=2,FakeSlant=0.2}]
\usepackage{titlesec}
\usepackage{indentfirst}

\usepackage[a4paper, 
            top=1cm, 
            bottom=2cm, 
            left=2.5cm, 
            right=2.5cm]{geometry}
\setcounter{secnumdepth}{0}

\titleformat{\subsection}
  {\normalfont\large\bfseries}{}{0pt}{}
\titleformat{\subsubsection}
  {\normalfont\large\bfseries}{}{0pt}{}
\titlespacing*{\subsubsection}
  {\parindent}{3.25ex plus 1ex minus .2ex}{1.5ex plus .2ex}
\newfontface\englabelfont{lmroman9-italic.otf}[
  Path=C:/texlive/2026/texmf-dist/fonts/opentype/public/lm/]
\newcommand{\englabel}[1]{{\normalfont\englabelfont\small #1}}

\usepackage{titling}
\usepackage{url}
\makeatletter
\renewcommand{\@makefntext}[1]{%
  \noindent
  \setlength{\hangindent}{2.133em}%
  \hangafter=1\relax
  \makebox[1.8em][r]{\@makefnmark}\hspace{0.333em}#1}
\makeatother
\title{\textbf{연구계획서 초안}} % chktex 36
\author{언어학과 이승헌(2025\mbox{-}17960)} % chktex 36
\date{}
\setlength{\droptitle}{-2em}
\pretitle{\begin{center}\LARGE}
\posttitle{\par\end{center}\vspace{-1em}}
\preauthor{\begin{center}\large}
\postauthor{\par\end{center}}
\predate{}
\postdate{}

\begin{document}
\maketitle

\subsection{연구 주제 \englabel{topic}}
수업 시간에 다룬 `영아의 말소리 지각 \textendash{} 단어 분절'\footnote{Jusczyk \& Aslin (1995) 강독}을 연구 주제로 선정하였다. 영아가 말소리를 지각하기 위해서는 크게 세 가지의 문제를 해결해야 하는데, 이 중 하나가 바로 단어 분절 문제다. 영아에게는 띄어쓰기 없는 말의 흐름에서 단어를 찾아야 하는 과제가 주어지며, Jusczyk \& Aslin (1995)에 따르면 생후 7.5개월 영아에게서 단어 분절 능력이 나타난다. 단, Jusczyk \& Aslin (1995)에서는 연령차에 따른 단어 분절 능력의 차이를 따로 검정하지 않았으므로, 단어 분절 능력이 어느 특정한 시점에 나타나는 것인지는 알 수 없다. 또한, 영어 습득 영아를 대상으로 실험을 진행했기 때문에, 단어 분절 능력 출현 시기에 대한 결과를 언어 보편적으로 확장할 수도 없다. 

\smallskip
\subsection{레퍼런스 \englabel{ref}} 
\subsubsection{출발 논문} 
\begin{itemize}
  \item Jusczyk, P. W., \& Aslin, R. N. (1995). Infants' detection of the sound patterns of words in fluent speech. \textit{Cognitive Psychology, 29}(1), 1--23.\ \url{https://doi.org/10.1006/cogp.1995.1010}
\end{itemize}

\vspace{-2ex}
\subsubsection{참고 문헌}
\begin{itemize}
    \item Chai, J. H., \& Ko, E.-S. (2025). Input properties shape word segmentation performance across child development: A computational modeling study. \textit{Infancy, 30}(5), e70046.\ \\ \url{https://doi.org/10.1111/infa.70046} % chktex 36
    \item Chai, J. H., Mun, S., \& Ko, E.-S. (2026). The alignment between language properties and computational algorithms enhances statistical word segmentation: Evidence from Korean child-directed speech. \textit{Journal of Child Language}. Advance online publication.\ \\ \url{https://doi.org/10.1017/S0305000926100646}
\end{itemize}

\smallskip
\subsection{연구 질문 \englabel{question}}
본 연구의 연구 질문은 `한국어 습득 6개월 및 7.5개월 영아는 한국어 발화에서 친숙한 단어 형태를 분절할 수 있는가? 단어 지각 양상은 영어 습득 영아에게서 보고된 결과와 유사한가?'이다. Jusczyk \& Aslin (1995)의 연구는 영어 습득 영아에게 영어 발화를 들려주었을 때의 단어 분절 양상만 관찰하였으므로, 다른 언어 환경에서도 유사한 실험 결과가 나타나는지는 별도의 실험을 통해 검증해 보아야 한다. 이에, 본 연구는 한국어 습득 영아를 대상으로 한국어 발화를 들려주었을 때의 단어 분절 양상을 관찰하고자 한다. % chktex 38 
선행 연구에서 한국어 습득 영아의 발달 단계에 따른 단어 분절 능력을 계산하긴 하였으나, 이는 영아의 실제 반응을 관찰한 결과가 아니라는 점에서 한계를 갖는다(Chai \& Ko, 2025; Chai et al., 2026).   

\smallskip
\subsection{가설 \englabel{hypothesis}}
\begin{itemize}
  \item 한국어 습득 7.5개월 영아는 발화에서 친숙한 단어 형태를 지각하지만, 6개월 영아에게서는 단어 지각이 뚜렷하게 나타나지 않을 것이다. 
\end{itemize}

\smallskip
\subsection{연구 방법 \englabel{methodology}}
Jusczyk \& Aslin (1995)의 연구 방법을 일부 참고할 계획이다. 실험은 `실험 1--실험 2', 총 2단계로 구성한다. 실험에서는 친숙화 단계를 통해 특정 단어를 반복하여 들려주고, 친숙화 과정에서 들은 단어가 포함된 문장과 그렇지 않은 문장을 들려준 뒤 문장별 청취 시간을 비교한다. 이때, 실험 1은 생후 7.5개월 한국어 습득 영아를 대상으로 하고, 실험 2는 생후 6개월 한국어 습득 영아를 대상으로 한다. 두 실험 결과를 비교함으로써 연령별 단어 지각 능력을 비교하고, 가설을 검증할 수 있을 것이다. 
\end{document}
